\section{Einseitige Nebenbedingungen}
Nun betrachten wir Variationsrechnungsprobleme mit einer punktweisen, einseitigen Nebenbedingung.
Wir definieren uns das Energiefunktional
\begin{equation}\label{equ:2.1}
I(w) :=
\int_U \frac{1}{2} \cdot | Dw |^2 - f  \cdot w \td{x}
\end{equation}
für alle Funktionen $w$, welche in der Menge
\begin{equation}\label{equ:2.2}
\mathcal{A}
:= 
\lbrace w \in H_0^1(U) \ | \ w \geq h \ \text{f.ü.} \rbrace
\end{equation}
enthalten sind. Außerdem sind die Funktionen $h : \overline{U} \to \R$ und $f : \R \to \R$ glatt.
Durch schnelles Nachrechnen erkennen wir, dass $\mathcal{A}$ konvex ist.

\begin{genericthm}{Existenz des Minimieres}\label{skript:2.1}
Sei $\mathcal{A}$ nichtleer.
Dann exisitiert ein eindeutiges $u \in \mathcal{A}$, so dass
\begin{align*}
I(u) := \min \limits_{w \in \mathcal{A}} I(w) 
\end{align*}
erfüllt ist.
\end{genericthm}

\begin{proof}
\begin{enumerate}

\item[\textbf{(1)}]
Die Existenz folgt sofort mit ähnlichen Ideen wie in \ref{skript:1.2}.
Sei $a_n$ eine Minimierungsfolge, dann existiert eine schwach konvergente Teilfolge $a_{n_k}$.
Diese konvergiert wiederum in $L^2(U)$ und die Folgenglieder sind in $\mathcal{A}$.
Damit folgt auch $u \in \mathcal{A}$.

\item[\textbf{(2)}]
Angenommen $u$ und $\tilde{u}$ sind verschieden und lösen das Minimierungsproblem.
Dann folgt mithilfe binomischer Formeln und
\begin{align*}
I\left( \frac{u + \tilde{u}}{2} \right)
&=
\int_U \frac{1}{2} \cdot \left| \frac{Du + D\tilde{u}}{2} \right|^2 - f \cdot \left( \frac{u + \tilde{u}}{2} \right) \td{x}\\
&= 
\int_U \frac{1}{8} \cdot ( 2 \cdot |Du|^2 + 2 \cdot |D \tilde{u} |^2 - | Du - D\tilde{u}|^2)  - f \cdot \left( \frac{u + \tilde{u}}{2} \right) \td{x}\\
&<
\frac{1}{2} \cdot \int_U \frac{1}{2} \cdot|Du|^2 - f \cdot u  \td{x} 
+ 
\frac{1}{2} \cdot \int_U \frac{1}{2} \cdot | D \tilde{u} |^2 - f \cdot \tilde{u} \td{x}\\
&=
\frac{1}{2} \cdot( I(u) + I(\tilde{u}))
\end{align*}
ein Widerspruch zur Minimaliät von $u$ und $\tilde{u}$.
\end{enumerate}
\end{proof}

\begin{genericthm}{Variationsungleichung}\label{skript:2.2}
Sei $u \in \mathcal{A}$ die eindeutige Lösung von
\begin{align*}
I(u) = \min \limits_{w \in \mathcal{A}} I(w).
\end{align*}
Dann gilt
\begin{equation}\label{equ:2.3}
\int_U f \cdot (w - u) \td{x} 
\leq 
\int_U Du \cdot D(w- u) \td{x}
\end{equation}
für alle $w \in \mathcal{A}$.
\end{genericthm}

\begin{proof}
Sei $w \in \mathcal{A}$ beliebig.
Wir wissen, dass $\mathcal{A}$ konvex ist.
Somit gilt $  u + \tau \cdot (w - u) \in \mathcal{A}$ für $0 \leq \tau \leq 1$.
Wir definieren
\begin{align*}
i(\tau) := I (u + \tau \cdot (w - u) ).
\end{align*}
Nun ist $u$ ein eindeutiges Minimum, womit  $i(0) \leq i ( \tau ) $ für $0 \leq \tau \leq 1 $ gilt.
Damit erhalten wir auch $i^\prime(0) \geq 0$.
Mit dieser Erkenntnis bewaffnet betrachten wir den Differenzenquotient
\begin{align*}
\frac{i(\tau) - i(0)}{\tau}
= ...
= \int_U Du \cdot D(w-u) + \frac{\tau}{2} \cdot | D(w-u)|^2 - f \cdot (w- u) \td{x}
\end{align*}
mit $0 < \tau \leq 1$ und erhalten für $\tau \rightarrow 0 $  durch
\begin{align*}
0  \leq i^\prime(0) = \int_U Du \cdot D(w-u) - f\cdot (w- u) \td{x}
\end{align*}
die gewünschte Aussage.
\end{proof}

\begin{genericdf}{Anwendung der Variationsungleichung}\label{skript:2.3}
Sei $u \in \mathcal{A}$ mit 
\begin{align*}
I(u) = \min \limits_{w \in \mathcal{A}} I(w).
\end{align*}
Wir setzen nun voraus, dass $\partial U $ glatt ist.
Nach Kinderlehrer-Stampacchia gilt dann $u \in W^{2,\infty}(U)$.
Damit folgt dann, dass
\begin{align*}
O := \lbrace x \in U \ | \ u(x) > h(x) \rbrace
\end{align*}
offen und
\begin{align*}
A := \lbrace x \in U \ | \ u(x) = h(x) \rbrace
\end{align*}
abgeschlossen ist.
Sei $v \in C_c^\infty(O)$ beliebig, aber fest.
Dann gilt $w := u + \tau \cdot v  \in \mathcal{A}$ für $|\tau|$ hinreichend klein.
Nun folgt
\begin{align*}
\int_O Du \cdot D(w - u ) - f(w-u) \td{x} 
= 
\tau \cdot \int_O Du \cdot Dv - fv  \td{x}
\geq 0 
\end{align*}
mit der Variationsungleichung.
Da $\tau$ postiv und negativ sein kann, erhalten wir  
\begin{align*}
\int_O -\Delta u \cdot v \td{x} = \int_O Du \cdot Dv \td{x}  = \int_O f \cdot v \td{x}.
\end{align*}
Somit ist $u$ eine schwache Lösung von
\begin{equation}\label{equ:2.4}
- \Delta u = f
\end{equation}
auf $O$ und die Regularitätstheorie für elliptische PDEs liefert uns $u \in C^\infty(O)$.
Sei $v \in C_c^\infty(U)$ beliebig mit $v \geq 0 $.
Dann gilt für $w := u + \tau \cdot v  \in \mathcal{A}$ mithilfe der Variationsungleichung und partieller Integration
\begin{align*}
\int_U Du \cdot Dv - fv  \td{x} = \int_U -\Delta u \cdot v - f \cdot v  \geq 0,
\end{align*}
wodurch 
\begin{equation}\label{equ:2.5}
-\Delta u \geq f 
\end{equation}
fast überall in $U$ gilt.
\end{genericdf}
